\chapter{Dados e Metodologia}
\label{cap:metodologia-linha1}

Este capítulo descreve os dados e o protocolo metodológico do estudo, dedicado à difusão musical em charts brasileiros. A organização parte do pressuposto de que os conceitos de SIR, grafos heterogêneos, redes neurais em grafos e grafos temporais já foram apresentados na fundamentação teórica. Assim, o foco aqui é operacional: quais dados são usados, como eles são transformados em um grafo heterogêneo temporal, quais modelos serão comparados e como a avaliação será conduzida.

%=====================================================

\section{Visão geral da linha experimental}

O estudo tem como objetivo comparar um baseline epidemiológico, baseado em SIR, com uma arquitetura neural em grafo heterogêneo temporal para previsão de difusão musical. O domínio empírico é composto por músicas presentes em charts brasileiros de streaming, especialmente Top 200 e Viral 50. A hipótese operacional é que a estrutura relacional entre músicas, artistas e gêneros acrescenta informação preditiva que não está disponível em modelos ajustados apenas sobre curvas agregadas.

O desenho experimental segue três etapas. Primeiro, reproduz-se o baseline SIR sobre séries de viralidade e sucesso, de forma alinhada ao trabalho de referência \citep{oliveira2025brasnam}. Segundo, constrói-se um grafo heterogêneo com nós de música, artista e gênero, preservando atributos e tempos de observação das relações. Terceiro, treina-se uma GNN heterogênea com componente temporal para prever valores futuros de \emph{rank-score} e comparar seu desempenho com o baseline.

%=====================================================

\section{Fontes de dados}

A base principal é o MGD+, que fornece metadados musicais e atributos acústicos associados a músicas e artistas \citep{seufitelli2023dsw}. Esses dados são combinados a séries de charts brasileiros do Spotify, em particular rankings diários de Top 200 e Viral 50. O Top 200 é usado como aproximação de sucesso ou popularidade consolidada, enquanto o Viral 50 representa sinais de crescimento e circulação viral.

O universo atual do estudo contém 6.526 músicas, 1.701 artistas e 530 gêneros. Esses elementos formam os três tipos de nós do grafo heterogêneo. As músicas carregam atributos acústicos, indicadores de popularidade e variáveis derivadas dos charts; os artistas carregam estatísticas agregadas de presença em charts e diversidade de gêneros; os gêneros são representados por embeddings aprendidos durante o treinamento.

A cobertura temporal considerada vai de 2017 a 2021. Essa janela permite separar treinamento, validação e teste de forma cronológica, evitando que informações futuras contaminem a avaliação. A separação temporal também é compatível com a formulação de snapshots semanais usada no grafo.

%=====================================================

\section{Pré-processamento das séries musicais}

O pré-processamento das séries segue a lógica do estudo SIR de referência para manter comparabilidade. Rankings diários são transformados em séries numéricas de \emph{rank-score}, separando a dimensão de viralidade, associada ao Viral 50, e a dimensão de sucesso, associada ao Top 200. Essa transformação converte posições de ranking em valores comparáveis entre músicas e datas. Para uma música em uma posição $rank$ de um chart com posição máxima $\textit{max\_rank}$ (200 no Top 200 e 50 no Viral 50), o \emph{rank-score} é definido por

\begin{equation}
\label{eq:rank-score}
\textit{rank\_score} = \textit{max\_rank} - rank + 1,
\end{equation}

de modo que a posição \#1 receba o valor máximo. Sobre essa série aplica-se uma média móvel de sete dias, seguida de normalização \emph{min-max} para o intervalo $[0,\,0{,}5]$; dias em que a música não aparece no chart recebem um valor mínimo de $0{,}001$, evitando problemas numéricos no ajuste do SIR. Cada música gera, portanto, duas séries de aproximadamente 1.895 dias: uma de viralidade e uma de sucesso.

Músicas sem atributos acústicos completos não são removidas do universo. Em vez disso, os valores ausentes são imputados com a mediana e acompanhados por uma variável binária indicando ausência acústica. Essa decisão preserva a cobertura do dataset e permite que o modelo aprenda se a ausência de atributos carrega algum sinal sistemático.

As variáveis contínuas usadas como atributos são padronizadas por escore-$z$. Essa normalização é aplicada a atributos como popularidade, streams, dias em chart e estatísticas agregadas de artistas. O objetivo é evitar que diferenças de escala dominem a passagem de mensagens na GNN.

%=====================================================

\section{Baseline SIR}

O baseline SIR modela a trajetória de cada música como um processo de contágio agregado. A curva observada de \emph{rank-score} é ajustada por meio dos parâmetros $\beta$ e $\gamma$, que representam, respectivamente, taxa de adoção e taxa de remoção. A razão $R_0 = \beta / \gamma$ resume o potencial de espalhamento da música segundo a leitura epidemiológica.

Esse baseline tem duas funções no projeto. A primeira é metodológica: reproduzir o resultado do estudo de referência e verificar se o pipeline local alcança erros compatíveis. A segunda é conceitual: oferecer uma comparação interpretável para avaliar se o aumento de complexidade trazido por GNNs heterogêneas se traduz em ganho empírico.

A métrica principal para essa etapa é RMSE, calculada separadamente para viralidade e sucesso. O critério mínimo do estudo é obter desempenho compatível com o paper de referência dentro de uma margem aceitável. Uma vez validado esse baseline, a comparação passa a incluir a predição genuína de valores futuros, e não apenas o ajuste retroativo da curva observada.

%=====================================================

\section{Construção do grafo heterogêneo}

O grafo proposto possui três tipos de nós: \texttt{music}, \texttt{artist} e \texttt{genre}. Essa escolha reflete a hipótese de que a difusão musical depende simultaneamente das propriedades da faixa, do histórico dos artistas associados e da proximidade cultural expressa por gêneros.

% FIGURA: diagrama esquemático do grafo heterogêneo C-A-G (music-artist-genre) com os quatro tipos de aresta (performs, has_genre, cotrajectory, cooccurs) e setas tipadas. (adiada)

As músicas são descritas por vetores de 15 dimensões. Esses vetores combinam nove atributos acústicos, popularidade, indicador de conteúdo explícito, tipo de música, total de streams, número de dias em chart e a flag de ausência acústica. Os artistas são descritos por quatro dimensões: número de hits, número de hits colaborativos, anos de presença em charts e número de gêneros associados. Os gêneros são representados por embeddings de 32 dimensões aprendidos durante o treinamento.

O grafo contém quatro tipos principais de arestas. A relação \texttt{performs} conecta artistas a músicas e representa autoria ou participação. A relação \texttt{has\_genre} conecta artistas a gêneros. A relação \texttt{cotrajectory} conecta músicas que compartilham presença temporal nos charts e possuem trajetórias semelhantes. A relação \texttt{cooccurs} conecta gêneros que aparecem conjuntamente em agregações do dataset.

Cada tipo de aresta possui atributos específicos. Em \texttt{performs}, são registrados papel do artista, posição na lista de artistas e primeira semana observada. Em \texttt{cotrajectory}, são registrados dias em conjunto, distância média de posição, tipo de chart e primeira semana observada. Em \texttt{cooccurs}, são registrados peso, popularidade média, streams médios e primeira semana observada. Essa estrutura permite que o modelo diferencie relações por semântica e por momento de observação.

%=====================================================

\section{Temporalidade e prevenção de vazamento}

A temporalidade é incorporada por meio do atributo \texttt{first\_seen\_week}. Para qualquer semana $t$, o modelo só pode acessar arestas observadas até aquela semana. Assim, snapshots do grafo são obtidos por filtragem temporal, sem materializar todas as versões em disco.

Essa escolha é importante porque relações como \texttt{cotrajectory} podem depender de observações futuras se forem calculadas de forma ingênua. Ao filtrar por primeira semana observada, a avaliação se aproxima do cenário real de previsão: em uma data $t$, o modelo deve prever o futuro usando apenas músicas, atributos e relações disponíveis até $t$.

A formulação por snapshots semanais foi escolhida por compatibilidade com o custo computacional e com o grão temporal dos dados processados. Modelos temporais contínuos, como TGN, permanecem como alternativa para a Linha 2, na qual eventos de adoção de trends em vídeo curto podem ser representados com timestamps mais finos.

%=====================================================

\section{Modelo proposto}

O modelo proposto combina uma GNN heterogênea com um componente recorrente temporal.

%=====================================================

\section{Protocolo experimental}

A avaliação usa splits temporais. O período de 2017 a junho de 2020 é usado para treinamento; julho a dezembro de 2020 é usado para validação; e 2021 é reservado para teste. Essa divisão impede que o modelo seja avaliado em músicas ou relações futuras vistas durante o treinamento.

Serão executados dois regimes de avaliação. O primeiro é o ajuste retroativo, voltado à comparação direta com o baseline SIR e com o trabalho de referência. Nesse modo, o objetivo é verificar se o modelo reproduz ou melhora a descrição das curvas já observadas. O segundo é a predição genuína, em que o modelo prevê \emph{rank-score} em horizontes futuros usando apenas dados disponíveis até o instante de previsão.

Os horizontes planejados são $k \in \{1, 7, 14, 30\}$ dias. As métricas principais são RMSE por horizonte e acerto direcional, isto é, se o modelo acerta a tendência de subida ou queda da música. Quando viável, também será considerada uma métrica probabilística, como CRPS, caso o modelo seja estendido para produzir distribuições preditivas.

%=====================================================

%=====================================================

\section{Síntese do capítulo}

A metodologia proposta transforma séries de charts musicais em um problema de previsão temporal sobre grafo heterogêneo. O baseline SIR oferece uma referência interpretável e já consolidada na literatura brasileira recente. O grafo heterogêneo adiciona relações entre músicas, artistas e gêneros, enquanto a arquitetura temporal busca explorar a evolução dessas relações sem vazamento de informação futura. Com isso, a avaliação poderá responder se a estrutura heterogênea temporal melhora a modelagem de difusão musical em relação a modelos epidemiológicos populacionais.